 %%%%%%%%%%%%%%%%%%%%%%%%%%%%%%%%%%%%%%%%%
% Beamer Presentation
% LaTeX Template
% Version 1.0 (10/11/12)
%
% This template has been downloaded from:
% http://www.LaTeXTemplates.com
%
% License:
% CC BY-NC-SA 3.0 (http://creativecommons.org/licenses/by-nc-sa/3.0/)
%
%%%%%%%%%%%%%%%%%%%%%%%%%%%%%%%%%%%%%%%%%

%----------------------------------------------------------------------------------------
%	PACKAGES AND THEMES
%----------------------------------------------------------------------------------------

\documentclass[9pt,aspectratio=169]{beamer}
\usepackage[utf8]{inputenc}
\usepackage{bm}
\usepackage{bbm}
\usepackage{graphicx}
\usepackage{cmbright}
\usefonttheme{professionalfonts}
\DeclareSymbolFont{greekletters}{OML}{sansmath}{m}{it}
\DeclareMathSymbol{\beta}{\mathord}{greekletters}{"0C}
%\usefonttheme{serif}
% The Beamer class comes with a number of default slide themes
% which change the colors and layouts of slides. Below this is a list
% of all the themes, uncomment each in turn to see what they look like.

%\usetheme{default}
%\usetheme{AnnArbor}
% \usetheme{Antibes}
%\usetheme{Bergen}
%\usetheme{Berkeley}
%\usetheme{Berlin}
%\usetheme{Boadilla}
% \usetheme{CambridgeUS}
%\usetheme{Copenhagen}
%\usetheme{Darmstadt}
% \usetheme{Dresden}
\usetheme{Frankfurt}
%\usetheme{Goettingen}
% \usetheme{Hannover}
% \usetheme{Ilmenau}
% \usetheme{JuanLesPins}
%\usetheme{Luebeck}
%\usetheme{Madrid}
%\usetheme{Malmoe}
%\usetheme{Marburg}
% \usetheme{Montpellier}
% \usetheme[width=2.5cm]{PaloAlto}
%\usetheme{Pittsburgh}
% \usetheme{Rochester}
%\usetheme{Singapore}
% \usetheme{Szeged}
% \usetheme{Warsaw}
% \setbeamertemplate{title page}[default][colsep=-4bp,rounded=false]
\setbeamertemplate{headline}{}
\defbeamertemplate{footline}{myframe number}
{
  \hspace{1pt}
    \usebeamercolor[structure]{page number in head/foot}%
  \usebeamerfont{page number in head/foot}%
  \raisebox{1.7ex}[0pt][0pt]{% <--- change here
    \insertframenumber\,/\,\inserttotalframenumber\kern1em}%
}
\setbeamertemplate{footline}[myframe number]
\AtBeginSection[]
{
  \begin{frame}
    \frametitle{Outline}
    \begin{minipage}{1.0\linewidth}
      \tableofcontents[currentsection,currentsubsection]
    \end{minipage}
  \end{frame}
}

\AtBeginSubsection[]
{
  \begin{frame}
    \frametitle{Outline}
    \begin{minipage}{1.0\linewidth}
      \tableofcontents[currentsection,currentsubsection]
    \end{minipage}
  \end{frame}
}

% As well as themes, the Beamer class has a number of color themes
% for any slide theme. Uncomment each of these in turn to see how it
% changes the colors of your current slide theme.
%\definecolor{mycolor}{rgb}{.8,.0,.3}
%\setbeamercolor*{palette primary}{use=structure,fg=white,bg=mycolor}
%\usecolortheme{albatross}
%\usecolortheme{beaver}
%\usecolortheme{beetle}
%\usecolortheme{crane}
%\usecolortheme{dolphin}
%\usecolortheme{dove}
%\usecolortheme{fly}
%\usecolortheme{lily}
%\usecolortheme{orchid}
% \usecolortheme{rose}
%\usecolortheme{seagull}
\usecolortheme{seahorse}
%\usecolortheme{whale}
%\usecolortheme{wolverine}
%\usecolortheme[named=mycolor]{structure}
%\setbeamertemplate{footline} % To remove the footer line in all slides uncomment this line
%\setbeamertemplate{footline}[page number] % To replace the footer line in all slides with a simple slide count uncomment this line

%\setbeamertemplate{navigation symbols}{} % To remove the navigation symbols from the bottom of all slides uncomment this line


\usepackage{booktabs} % Allows the use of \toprule, \midrule and \bottomrule in tables

% Custom commands
\newcommand{\RR}{\mathbb{R}}
\newcommand{\ZZ}{\mathbb{Z}}
\newcommand{\X}{\mathcal{X}}
\newcommand{\Y}{\mathcal{Y}}

\title{Lecture 9: Introduction to the Theory of Optimal Control}

\author{Junnan Zhang} % Your name
\institute[Paula and Gregory Chow Institute for Studies in Economics\\
Xiamen University] % Your institution as it will appear on the bottom of every slide, may be shorthand to save space
{
	Paula and Gregory Chow Institute for Studies in Economics \\ Xiamen University  \\ % Your institution for the title page
	\medskip
}
\date{Fall, 2026} % Date, can be changed to a custom date
\setbeamertemplate{itemize items}[default]
\setbeamertemplate{sidebar}[sidebar right]
% \definecolor{myco}{HTML}{a67678}
\definecolor{myco}{HTML}{8776a6}
\setbeamercolor{itemize item}{fg=myco} % all frames will have red bullets
\setbeamercolor{itemize subitem}{fg=myco} % all frames will have red bullets
\setbeamercolor{block title}{bg=myco}
\setbeamercolor{structure}{fg=myco}
% \setbeamercovered{transparent}


\begin{document}

\frame{\titlepage}

\begin{frame}{Introduction to Continuous-Time Optimization}
    \begin{itemize}
        \item This chapter presents a number of basic results in dynamic
        optimization in continuous time, particularly the so-called optimal
        control approach.
        \item Both dynamic optimization in discrete time and in continuous
        time are useful tools for macroeconomics and other areas of dynamic
        economic analysis.
        \item One approach is not superior to the other; instead, certain
        problems become simpler in discrete time while others are naturally
        formulated in continuous time.
    \end{itemize}
\end{frame}

\begin{frame}{Mathematical Challenges in Continuous-Time Optimization}
    \begin{itemize}
        \item Continuous-time optimization introduces several new
        mathematical issues.
        \item This is largely because even with a finite horizon, the
        maximization is with respect to an infinite-dimensional object: we
        are maximizing over an entire function: $y : [t_0, t_1] \rightarrow
        \mathbb{R}$.
    \end{itemize}
\end{frame}

\begin{frame}{The Canonical Continuous-Time Optimization Problem}
The canonical continuous-time optimization problem can be written as:
\begin{align*}
\max_{x(t),y(t)} W(x(t), y(t)) &\equiv \int_0^{t_1} f(t, x(t), y(t))dt
\end{align*}
subject to:
\begin{align*}
\dot{x}(t) &= G(t, x(t), y(t))\\
x(t) &\in \X(t), \quad y(t) \in \Y(t) \text{ for all } t\\
x(0) &= x_0
\end{align*}
where for each $t$, $x(t)$ and $y(t)$ are finite-dimensional vectors.

\begin{itemize}
\item For each $t$: $\X(t) \subset \mathbb{R}^{K_x}$ and $\Y(t) \subset \mathbb{R}^{K_y}$, where $K_x, K_y \in \mathbb{N}$.
\item The vector $x$ denotes the \textbf{state variables}, which is governed by a
system of differential equations, given the behavior of the vector of \textbf{control
variables} $y$.
\item The end of the planning horizon $t_1$ can be equal to infinity.
\end{itemize}
\end{frame}

\section{Variational Arguments}

\begin{frame}{Special Case: One-Dimensional Problem}
Consider the following special case where the horizon is finite and both the state and control variables are one dimensional:
\begin{align*}
\max_{x(t),y(t),x_1} W(x(t), y(t)) &\equiv \int_0^{t_1} f(t, x(t), y(t))dt
\end{align*}
subject to:
\begin{align*}
\dot{x}(t) &= g(t, x(t), y(t))\\
x(t) &\in \X, \quad y(t) \in \Y \text{ for all } t\\
x(0) &= x_0, \quad x(t_1) = x_1
\end{align*}
\begin{itemize}
    \item The sets $\X(t)$ and $\Y(t)$ in the more general problem are now
    taken to be independent of time for simplicity. We assume that
    $\X$ and $\Y$ are nonempty and convex.
    \item Notice that there is also a terminal value constraint $x(t_1) =
    x_1$, but $x_1$ is included as an additional choice variable.
\end{itemize}
\end{frame}

\begin{frame}{Admissible Pairs and Assumptions}
\begin{itemize}
\item A pair of functions $(x(t), y(t))$ that satisfies the constraint and boundary conditions is referred to as an \textbf{admissible pair}.
\item Throughout, as in the previous chapter, we suppose that the value of
the objective function is finite. That is, $W(x(t), y(t)) < \infty$ for any
admissible pair $(x(t), y(t))$.
\item Let us first suppose that $t_1 < \infty$, so that we have a
finite-horizon optimization problem.
\item Assume that $f$ and $g$ are continuously differentiable.
\end{itemize}
\end{frame}

\begin{frame}{Key Challenges in Characterizing Solutions}
The challenge in characterizing the optimal solution to this problem lies in two features:
\begin{enumerate}
\item We are choosing a function $y : [0, t_1] \rightarrow \Y$ rather than a vector or a finite-dimensional object.
\item The constraint takes the form of a differential equation rather than a set of inequalities or equalities.
\end{enumerate}

These features make it difficult for us to know what type of optimal policy to look for:
\begin{itemize}
\item $y$ may be a highly discontinuous function.
\item It may also hit the boundary of the feasible set, thus corresponding to
a corner solution.
\end{itemize}
\end{frame}

\begin{frame}{The Variational Approach}
Fortunately, in most economic problems
\begin{itemize}
    \item There is enough structure to make solutions continuous functions.
    \item The Inada conditions ensure that solutions lie in the interior of
    the feasible set.
\end{itemize}
Then, it can be characterized by using variational arguments:
\begin{itemize}
    \item First assume that there exists a continuous solution (function)
    $\hat{y}$ that lies everywhere in the interior of the set $\Y$.
    \item Then give necessary conditions
\end{itemize}
\end{frame}

\begin{frame}{Formal Setup for Variational Approach}
More formally, let us assume that $(\hat{x}(t), \hat{y}(t))$ is an admissible pair such that:
\begin{itemize}
\item $\hat{y}(\cdot)$ is continuous on $[0, t_1]$
\item $(\hat{x}(t), \hat{y}(t)) \in \text{Int } \X \times \text{Int } \Y$
\end{itemize}

And that we have:
$$W(\hat{x}(t), \hat{y}(t)) \geq W(x(t), y(t))$$
for any other admissible pair $(x(t), y(t))$.
\begin{itemize}
    \item Continuous and interior solution is a strong assumption
    \item When $y(t)$ is continuous, $\dot{x}(t)$ will also be continuous,
    and $x(t)$ is continuously differentiable.
\end{itemize}
\end{frame}

\begin{frame}{Constructing Variations}
    \begin{itemize}
        \item Take an arbitrary fixed continuous function $\eta(t)$ and let
        $\varepsilon \in \mathbb{R}$ be a real number. Then a variation of
        the function $\hat{y}(t)$ is defined by $y(t, \varepsilon) \equiv
        \hat{y}(t) + \varepsilon\eta(t)$.
        \item Let us also define $x(t, \varepsilon)$ as the path of the state
        variable corresponding to the path of control variable $y(t,
        \varepsilon)$ with $x(0, \varepsilon) = x_0$.
        \item Since $(\hat{x}(t), \hat{y}(t))$ is interior and continuous on
        a compact set, we can always find $\varepsilon \in
        [-\varepsilon_\eta, \varepsilon_\eta]$ such that $(x(t, \varepsilon),
        y(t, \varepsilon))$ is an admissible pair.
    \end{itemize}
\end{frame}

\begin{frame}{Optimality Condition}
    \begin{itemize}
        \item Define:
    \begin{align*}
        W(\varepsilon) \equiv W(x(t, \varepsilon), y(t, \varepsilon)) =
        \int_0^{t_1} f(t, x(t, \varepsilon), y(t, \varepsilon))dt 
    \end{align*}
    \item Since $\hat{y}(t)$ is optimal , we have:
    $W(\varepsilon) \leq W(0) \text{ for all } \varepsilon \in
    [-\varepsilon_\eta, \varepsilon_\eta]$
    \item Rewrite the differential equation constraint: $g(t, x(t,
    \varepsilon), y(t, \varepsilon)) - \dot{x}(t, \varepsilon) = 0$ for all
    $t \in [0, t_1]$.
    \item Thus for any function $\lambda : [0, t_1] \rightarrow \mathbb{R}$,
    we have:
    \begin{equation*}
        \int_0^{t_1} \lambda(t)[g(t, x(t, \varepsilon), y(t, \varepsilon)) -
        \dot{x}(t, \varepsilon)] dt = 0
    \end{equation*}
    \item Adding the constraint equation to the objective function yields:
    \begin{equation*}
        W(\varepsilon) = \int_0^{t_1} \{f(t, x(t, \varepsilon), y(t,
        \varepsilon)) + \lambda(t)[g(t, x(t, \varepsilon), y(t, \varepsilon))
        - \dot{x}(t, \varepsilon)]\}dt
    \end{equation*}
\end{itemize}

\end{frame}

\begin{frame}{First-Order Necessary Condition}
    \begin{itemize}
        \item We can integrate $\int_0^{t_1}\lambda(t) \dot x(t, \varepsilon)
        dt$ by part to obtain
        \begin{equation*}
            \int_0^{t_1}\lambda(t) \dot x(t, \varepsilon) dt = \lambda(t_1)
            x(t_1, \varepsilon) - \lambda(0) x_0 - \int_0^{t_1}\dot
            \lambda(t) x(t, \varepsilon) dt
        \end{equation*}
        \item Substituting this back to $W(\varepsilon)$ and differentiating
        $W(\varepsilon)$ with respect to $\varepsilon$ gives:
    \begin{align*}
        W'(\varepsilon) \equiv &\int_0^{t_1} [f_x(t, x(t, \varepsilon), y(t,
        \varepsilon)) + \lambda(t)g_x(t, x(t, \varepsilon), y(t,
        \varepsilon)) + \dot{\lambda}(t)]x_\varepsilon(t, \varepsilon)dt\\ 
        &+\int_0^{t_1} [f_y(t, x(t, \varepsilon), y(t, \varepsilon)) +
        \lambda(t)g_y(t, x(t, \varepsilon), y(t, \varepsilon))]\eta(t)dt -
        \lambda(t_1)x_\varepsilon(t_1, \varepsilon)
    \end{align*}
    \item Consequently, optimality requires that:
    $$W'(0) = 0 \textbf{ for all } \eta(t)$$
\end{itemize}
\end{frame}

\begin{frame}{First-Order Necessary Condition}
    \begin{itemize}
        \item Evaluate $W'(\varepsilon)$ at $\varepsilon = 0$ to obtain:
        \begin{align*}
            W'(0) \equiv &\int_0^{t_1} [f_x(t, \hat{x}(t), \hat{y}(t)) +
            \lambda(t)g_x(t, \hat{x}(t), \hat{y}(t)) +
            \dot{\lambda}(t)]x_\varepsilon(t, 0)dt\\ 
            &+\int_0^{t_1} [f_y(t, \hat{x}(t), \hat{y}(t)) + \lambda(t)g_y(t,
            \hat{x}(t), \hat{y}(t))]\eta(t)dt -
            \lambda(t_1)x_\varepsilon(t_1, 0)
        \end{align*}
        \item Since it applies for any continuously differentiable
        $\lambda(t)$ function, let us consider the function $\lambda(t)$ that
        is a solution to the differential equation:
        $$\dot{\lambda}(t) = -[f_x(t, \hat{x}(t), \hat{y}(t)) + \lambda(t)g_x(t,
        \hat{x}(t), \hat{y}(t))]$$ with boundary condition $\lambda(t_1) =
        0$. This equation has a solution when $f_x$ and $g_x$ are continuous.
        \item Defining $\lambda(t)$ in this way allows us to isolate the
        effect of variations in the control variable $y(t)$
    \end{itemize}
\end{frame}

\begin{frame}{First-Order Necessary Condition}
    \begin{itemize}
        \item Since $\eta(t)$ is arbitrary, $\lambda(t)$ and $(\hat{x}(t),
        \hat{y}(t))$ need to be such that:
        \begin{equation*}
            f_y(t, \hat{x}(t), \hat{y}(t)) + \lambda(t)g_y(t, \hat{x}(t),
            \hat{y}(t)) = 0 \text{ for all } t \in [0, t_1]
        \end{equation*}
        \item Otherwise, there would exist some $\eta(t)$ that would make the
        following integral nonzero, contradicting optimality
        \begin{equation*}
            \int_0^{t_1} [f_y(t, \hat{x}(t), \hat{y}(t)) + \lambda(t)g_y(t,
            \hat{x}(t), \hat{y}(t))]\eta(t)dt
        \end{equation*}
        \item This argument establishes the necessary conditions for
        $(\hat{x}(t), \hat{y}(t))$ to be an interior continuous solution
    \end{itemize}
\end{frame}

\begin{frame}{Necessary Conditions: Free Terminal Value}
    \begin{block}{Theorem 7.1 (Necessary Conditions)}
        Consider the problem of maximizing
        \begin{equation*}
            \max_{x(t),y(t),x_1} W(x(t), y(t)) \equiv \int_0^{t_1} f(t, x(t), y(t))dt
        \end{equation*}
        subject to $\dot{x}(t) = g(t, x(t), y(t))$ and $x(t) \in \X$, $y(t)
        \in \Y$ for all $t$, $x(0) = x_0$ and $x(t_1) = x_1$, with $f$ and
        $g$ continuously differentiable. Suppose that this problem has an
        interior continuous solution $(\hat{x}(t), \hat{y}(t)) \in \text{Int
        } \X \times \text{Int } \Y$. Then there exists a continuously differentiable
        costate function $\lambda(\cdot)$ defined on $t \in [0, t_1]$ such
        that the following are satisfied
        \begin{enumerate}
            \item $\dot{x}(t) = g(t, x(t), y(t))$
            \item $\dot{\lambda}(t) = -[f_x(t, \hat{x}(t), \hat{y}(t)) +
            \lambda(t)g_x(t, \hat{x}(t), \hat{y}(t))]$
            \item $f_y(t, \hat{x}(t), \hat{y}(t)) + \lambda(t)g_y(t,
            \hat{x}(t), \hat{y}(t)) = 0$
            \item $\lambda(t_1) = 0$
        \end{enumerate}
    \end{block}
\end{frame}

\begin{frame}{The Transversality Condition}
\begin{itemize}
    \item The condition that $\lambda(t_1) = 0$ is the \textbf{transversality
      condition} of continuous-time optimization problems.
    \item It is naturally related to the transversality condition we
    encountered in the discrete-time case.
    \item Intuitively, this condition captures the fact that after the
    planning horizon, there is no value to having more (or less) $x$.
\end{itemize}
\end{frame}

\begin{frame}{Alternative Formulation: Fixed Terminal Value}
    \begin{block}{Theorem 7.2 (Necessary Conditions II)}
        Consider the problem of maximizing
        \begin{equation*}
            \max_{x(t),y(t)} W(x(t), y(t)) \equiv \int_0^{t_1} f(t, x(t), y(t))dt
        \end{equation*}
        subject to $\dot{x}(t) = g(t, x(t), y(t))$ and $x(t) \in \X$, $y(t)
        \in \Y$ for all $t$, $x(0) = x_0$ and $x(t_1) = x_1$, with $f$ and
        $g$ continuously differentiable. Suppose that this problem has an
        interior continuous solution $(\hat{x}(t), \hat{y}(t)) \in \text{Int
        } \X \times \text{Int } \Y$. Then there exists a continuously differentiable
        costate function $\lambda(\cdot)$ defined over $t \in [0, t_1]$ such
        that the following are satisfied:
        \begin{enumerate}
            \item $\dot{x}(t) = g(t, x(t), y(t))$
            \item $\dot{\lambda}(t) = -[f_x(t, \hat{x}(t), \hat{y}(t)) +
            \lambda(t)g_x(t, \hat{x}(t), \hat{y}(t))]$
            \item $f_y(t, \hat{x}(t), \hat{y}(t)) + \lambda(t)g_y(t,
            \hat{x}(t), \hat{y}(t)) = 0$
        \end{enumerate}
    \end{block}
    The transversality condition $\lambda(t_1) = 0$ is no longer present, but
    instead the terminal value of the state variable $x$ is specified as part
    of the constraints.
\end{frame}

\begin{frame}{Example 7.1}
    Consider the utility-maximizing problem of a consumer living between two
    dates, 0 and 1:
    \begin{align*}
        \max_{[c(t),a(t)]_{t=0}^1} &\int_0^1 \exp(-\rho t)u(c(t))dt\\
        \text{s.t. } \dot{a}(t) &= ra(t) + w - c(t)\\
        a(t) &\geq 0
    \end{align*}
    with the initial value of $a(0) > 0$. In this problem:
    \begin{itemize}
        \item Assume $u : \mathbb{R}_+ \rightarrow \mathbb{R}$ is strictly
        increasing, continuously differentiable, and strictly concave
        \item Consumption is the control variable
        \item The asset holdings of the individual are the state variable
    \end{itemize}
\end{frame}

\begin{frame}{Example 7.1}
\begin{itemize}
    \item To be able to apply Theorem 7.2, we need a terminal condition for
    $a(t)$. The economics of the problem implies that $a(1) = 0$.
    \item Theorem 7.2 provides the following necessary conditions for an
    interior continuous solution: there exists a continuously differentiable
    costate variable $\lambda(t)$ that satisfies
    \begin{itemize}
        \item a consumption Euler equation
        \begin{equation}
            \exp(-\rho t)u'(\hat{c}(t)) = \lambda(t)\tag{7.14}            
        \end{equation}
        \item a differential equation
        \begin{equation}
            \dot{\lambda}(t) = -r\lambda(t) \tag{7.15}
        \end{equation}
    \end{itemize}
    \item Using (7.15) and differentiating the first-order condition (7.14)
    yields a differential equation in consumption.
    \item We can also integrate (7.15) to obtain $\lambda(t) = \lambda(0)
    \exp(-rt)$. Combining this equation with (7.14) yields the optimal
    consumption: $\hat{c}(t) = u'^{-1}[\lambda(0) \exp((\rho - r)t)]$
\end{itemize}
\end{frame}

\begin{frame}{Example 7.1: Consumption Patterns}
    This equation implies different consumption patterns depending on the
    relationship between $\rho$ and $r$:
    \begin{itemize}
        \item When $\rho = r$, so that the discount factor and the rate of
        return on assets are equal, the individual will have a constant
        consumption profile.
        \item When $\rho > r$, the fact that $u'^{-1}$ is decreasing over
        time implies that consumption must be declining: a front-loaded
        consumption profile.
        \item When $\rho < r$, the opposite reasoning applies, and she
        chooses a back-loaded consumption profile.
    \end{itemize}
\end{frame}

\begin{frame}{Example 7.1: Determining Initial Consumption}
\begin{itemize}
\item The only variable left to determine to completely characterize the consumption profile is the initial value of the costate variable (and thus the initial value of consumption).
\item This comes from the observation that the individual will run down all her assets by the end of her planning horizon, that is, $a(1) = 0$.
\item Using the consumption rule, we have:
$\dot{a}(t) = ra(t) + w - u'^{-1}[\lambda(0) \exp((\rho - r)t)]$
\item The initial value of the costate variable, $\lambda(0)$, then has to be chosen such that $a(1) = 0$.
\end{itemize}
\end{frame}

\begin{frame}{Example 7.1: Applying Theorem 7.1?}
\begin{itemize}
    \item It may at first appear that Theorem 7.1 is more convenient to use
    than Theorem 7.2.
    \item The first-order necessary conditions still give: $\lambda(t) =
    \lambda(0) \exp(-rt)$. However, since $\lambda(1) = 0$, this equation
    holds only if $\lambda(t) = 0$ for all $t \in [0, 1]$.
    \item But $\exp(-\rho t)u'(\hat{c}(t)) = \lambda(t)$, which cannot be
    satisfied since $u' > 0$.
    \item Theorem 7.1 cannot be applied to this problem, because there is an
    additional constraint that $a(t) \geq 0$.
\end{itemize}
\end{frame}

\begin{frame}{Inequality Terminal Constraints}
    \begin{block}{Theorem 7.3 (Necessary Conditions III)}
        Consider the problem of maximizing
        \begin{equation*}
            \max_{x(t),y(t)} W(x(t), y(t)) \equiv \int_0^{t_1} f(t, x(t),
            y(t))dt
        \end{equation*}
        subject to $\dot{x}(t) = g(t, x(t), y(t))$, $(x(t), y(t)) \in \X
        \times \Y$ for all $t$, $x(0) = x_0$, and $x(t_1) \geq x_1$, with $f$
        and $g$ continuously differentiable. 

        Suppose that this problem has an interior continuous solution
        $(\hat{x}(t), \hat{y}(t)) \in \text{Int } \X \times \text{Int } \Y$.
        Then there exists a continuously differentiable costate function
        $\lambda(\cdot)$ defined over $t \in [0, t_1]$ such that the
        following are satisfied
        \begin{enumerate}
            \item $\dot{x}(t) = g(t, x(t), y(t))$
            \item $\dot{\lambda}(t) = -[f_x(t, \hat{x}(t), \hat{y}(t)) +
            \lambda(t)g_x(t, \hat{x}(t), \hat{y}(t))]$
            \item $f_y(t, \hat{x}(t), \hat{y}(t)) + \lambda(t)g_y(t,
            \hat{x}(t), \hat{y}(t)) = 0$
            \item $\lambda(t_1) \geq 0$
            \item $\lambda(t_1)(x(t_1) - x_1) = 0$ \quad
            (\emph{Complementary slackness condition})
        \end{enumerate}
    \end{block}
\end{frame}

\section{The Maximum Principle: A First Look}

\begin{frame}{The Hamiltonian Function}
\begin{itemize}
    \item By analogy with the Lagrangian, we can express the results more economically by constructing the Hamiltonian:
    $$H(t, x(t), y(t), \lambda(t)) \equiv f(t, x(t), y(t)) + \lambda(t)g(t, x(t), y(t))$$
    \item We often write $H(t, x, y, \lambda)$ for the Hamiltonian to simplify notation
    \item Since $f$ and $g$ are continuously differentiable, so is $H$
    \item We denote the partial derivatives of the Hamiltonian with respect to $x(t)$, $y(t)$, and $\lambda(t)$ by $H_x$, $H_y$, and $H_\lambda$, respectively
\end{itemize}
\end{frame}

\begin{frame}{Theorem 7.4: Simplified Maximum Principle (Part 1)}
\begin{block}{Theorem 7.4 (Simplified Maximum Principle)}
    Consider the problem of maximizing 
    \begin{equation*}
        \max_{x(t),y(t),x_1} W(x(t), y(t)) \equiv \int_0^{t_1} f(t, x(t), y(t))dt    
    \end{equation*}
    subject to $\dot{x}(t) = g(t, x(t), y(t))$ and $x(t) \in \X$, $y(t) \in
    \Y$ for all $t$, $x(0) = x_0$ and $x(t_1) = x_1$, with $f$ and $g$
    continuously differentiable. Suppose that this problem has an interior
    continuous solution $(\hat{x}(t), \hat{y}(t)) \in \text{Int } \X \times
    \text{Int } \Y$. Then there exists a continuously differentiable function
    $\lambda(t)$ such that the optimal control $\hat{y}(t)$ and the
    corresponding path of the state variable $\hat{x}(t)$ satisfy the
    following necessary conditions:
    \begin{align*}
        H_y(t, \hat{x}(t), \hat{y}(t), \lambda(t)) = 0\\
        \dot{\lambda}(t) = -H_x(t, \hat{x}(t), \hat{y}(t), \lambda(t)) \\ 
        \dot{x}(t) = H_\lambda(t, \hat{x}(t), \hat{y}(t), \lambda(t))
    \end{align*}
    for all $t \in [0, t_1]$ with $x(0) = x_0$ and $\lambda(t_1) = 0$.
    Moreover, the Hamiltonian $H(t, x, y, \lambda)$ also satisfies the
    Maximum Principle that $H(t, \hat{x}(t), \hat{y}(t), \lambda(t)) \geq
    H(t, \hat{x}(t), y, \lambda(t))$ for all $y \in \Y$ for all $t \in [0,
    t_1]$.
\end{block}
\end{frame}

\begin{frame}{Key Features of the Maximum Principle}
\begin{enumerate}
    \item As in the usual constrained maximization problems, a solution is
    characterized jointly with a set of "multipliers" $\lambda(t)$, and the
    optimal path of the control and state variables, $\hat{y}(t)$ and
    $\hat{x}(t)$
    
    \item The costate variable $\lambda(t)$ is informative about the value of
    relaxing the constraint (at time $t$). In particular, $\lambda(t)$ is the
    value of an infinitesimal increase in $x(t)$ at time $t$
    
    \item With this interpretation, it makes sense that $\lambda(t_1) = 0$ is
    part of the necessary conditions. After the planning horizon, there is no
    value to having more (or less) $x$. This is therefore the finite-horizon
    equivalent of the transversality condition in the previous chapter
\end{enumerate}
\end{frame}

\begin{frame}{Limitations of Necessary Conditions}
\begin{itemize}
    \item Theorem 7.4 gives necessary conditions for an interior continuous
    solution. However, we do not know whether such a solution exists
    \item Moreover, these necessary conditions may characterize a stationary
    point rather than a maximum or simply a local rather than a global
    maximum
    \item Sufficiency is again guaranteed by imposing concavity
\end{itemize}
\end{frame}

\begin{frame}{Mangasarian's Sufficiency Conditions}

\begin{block}{Theorem 7.5 (Mangasarian's Sufficiency Conditions)}
    Consider the same problem as above. Suppose that an interior continuous
    pair $(\hat{x}(t), \hat{y}(t)) \in \text{Int } \X \times \text{Int } \Y$ exists and
    satisfies the necessary conditions from Theorem 7.4. Suppose also that
    $\X \times \Y$ is a convex set and given the resulting costate variable
    $\lambda(t)$, $H(t, x, y, \lambda)$ is jointly concave in $(x, y) \in \X
    \times \Y$ for all $t \in [0, t_1]$. Then the pair $(\hat{x}(t),
    \hat{y}(t))$ achieves the global maximum of the objective function.
    Moreover, if $H(t, x, y, \lambda)$ is strictly concave in $(x, y)$ for
    all $t \in [0, t_1]$, then the pair $(\hat{x}(t), \hat{y}(t))$ is the
    unique solution.
\end{block}
\end{frame}

\begin{frame}{Arrow's Sufficiency Conditions}
    Define the maximized Hamiltonian:
    $$M(t, x(t), \lambda(t)) \equiv \max_{y \in \Y} H(t, x(t), y, \lambda(t))$$
\begin{block}{Theorem 7.6 (Arrow's Sufficiency Conditions)}
    Consider the problem as above and suppose that an interior continuous
    pair $(\hat{x}(t), \hat{y}(t)) \in \text{Int } \X \times \text{Int } \Y$ exists and
    satisfies the necessary conditions. Given the resulting costate variable
    $\lambda(t)$, if $\X$ is a convex set and $M(t, x, \lambda)$ is concave
    in $x \in \X$ for all $t \in [0, t_1]$, then $(\hat{x}(t), \hat{y}(t))$
    achieves the global maximum of the objective function. Moreover, if $M(t,
    x, \lambda)$ is strictly concave in $x$ for all $t \in [0, t_1]$, then
    the pair $(\hat{x}(t), \hat{y}(t))$ is the unique solution.
\end{block}

Theorem 7.6 weakens the concavity condition in Theorem 7.5 that $H(t, x, y,
\lambda)$ is jointly concave in $(x, y)$.
\end{frame}

\begin{frame}{Implications of Sufficiency Results}
\begin{itemize}
    \item One difficulty is verifying the concavity conditions in Theorem 7.5
    and Theorem 7.6
    \item Nevertheless, in many economically interesting situations, we can
    ascertain that the costate variable $\lambda(t)$ is everywhere
    nonnegative, for example, when $f_y(t, \hat{x}(t), \hat{y}(t)) > 0$
    and $g_y(t, \hat{x}(t), \hat{y}(t)) < 0$
    \item Once we know that $\lambda(t)$ is nonnegative, $H = f + \lambda g$
    is concave if $f$ and $g$ are both concave
\end{itemize}
\end{frame}

\begin{frame}
    \frametitle{Example}
    Solve the problem:
    \begin{align*}
        \max &\int_0^T \left[1 - tx(t) - u^2(t)\right]dt\\
        \text{s.t. } & \dot x(t) = u(t),\, x(0) = x_0 > 0,\, u(t) \in \RR
    \end{align*}
    \begin{itemize}
        \item The Hamiltonian is $H(t, x, u, \lambda) = 1 - tx - u^2 +
        \lambda u$, which is concave in $(x, u)$
        \item By the Maximum Principle, the following necessary conditions
        are satisfies:
        \begin{align*}
            H_u = -2\hat u(t) + \lambda(t) = 0\\
            \dot \lambda(t) = - H_x = t,\, \lambda(T) = 0\\
            \dot x(t) = u(t),\, x(0) = x_0
        \end{align*}
        \item Solution satisfies: $\lambda(t) = t^2/2 - T^2/2$, $\hat u(t) =
        t^2/4 - T^2/4$, and $\hat x(t) = t^3/12 - T^2t/4 + x_0$
    \end{itemize}
\end{frame}

\section{Infinite-Horizon Optimal Control}

\begin{frame}{Infinite-Horizon Problems}
\begin{itemize}
    \item The results presented so far are most useful in developing an
    intuition for how dynamic optimization in continuous time works.
    \item Most economic problems---including almost all growth models---are
    more naturally formulated as infinite-horizon problems.
    \item In this section, we provide necessary and sufficient conditions for
    optimality in infinite-horizon optimal control problems.
\end{itemize}
\end{frame}

\begin{frame}{Generalize the Terminal Value Constraint}
\begin{itemize}
    \item Throughout this chapter, let $b : \mathbb{R}_+ \to \mathbb{R}_+$,
    such that $\lim_{t \to \infty} b(t)$ exists and is finite.
    \item The terminal value condition is $\lim_{t \to \infty} b(t)x(t) \geq
    x_1$ for some $x_1 \in \mathbb{R}$.
    \item The special case where $b(t) \equiv 1$ gives us the terminal value
    constraint as $\lim_{t \to \infty} x(t) \geq x_1$ and is sufficient in
    many applications.
\end{itemize}
\end{frame}

\begin{frame}{The Infinite-Horizon Optimal Control Problem}
Using the same notation as above, the infinite-horizon optimal control problem is:
\begin{align}
    \max_{x(t),y(t)} W(x(t), y(t)) &\equiv \int_0^{\infty} f(t, x(t), y(t))dt \tag{7.32}\\
    \text{subject to }\quad \dot{x}(t) &= g(t, x(t), y(t)), \tag{7.33}\\
    x(t)\in \X,\, y(t) \in \Y \text{ for all } t,\, x(0) &= x_0 \text{ and }
    \lim_{t \to \infty} b(t)x(t) \geq x_1, \tag{7.34}
\end{align}
\end{frame}

\begin{frame}{Key Differences from Finite-Horizon Case}
\begin{itemize}
    \item The main difference is that now time runs to infinity.
    \item This problem allows for an implicit choice over the endpoint $x_1$,
    since there is no terminal date. The last part of (7.34) imposes a lower
    bound on this endpoint.
    \item $\X$ and $\Y$ need not be bounded sets.
    \item An admissible pair $(x(t), y(t))$ is defined in the same way as
    above, except that $y(t)$ can now be a piecewise continuous function.
\end{itemize}
\end{frame}

\begin{frame}{The Value Function}
    Define the value function, which is the analogue of the value function in
    discrete-time dynamic programming introduced in the previous chapter:
    \begin{align}
        V(t_0, x(t_0)) = \sup_{(x(t),y(t)) \in \X \times \Y}
        \int_{t_0}^{\infty} f(t, x(t), y(t))dt \tag{7.35}
    \end{align}
    subject to $\dot{x}(t) = g(t, x(t), y(t))$ and $\lim_{t \to \infty} b(t)x(t) \geq x_1$.
    \begin{itemize}
        \item $V(t_0, x(t_0))$ gives the optimal value starting at time $t_0$
        with state variable $x(t_0)$.
        \item $V(t_0, x(t_0)) \geq \int_{t_0}^{\infty} f(t, x(t), y(t))dt$
        for any admissible pair $(x(t), y(t))$.
        \item When $(\hat{x}(t), \hat{y}(t))$ is optimal, then $V(t_0,
        x(t_0)) = \int_{t_0}^{\infty} f(t, \hat{x}(t), \hat{y}(t))dt$.
    \end{itemize}
\end{frame}

\begin{frame}{Principle of Optimality}
    \begin{block}{Lemma 7.1 (Principle of Optimality)}
        Suppose that the pair $(\hat{x}(t), \hat{y}(t))$ is a solution to
        (7.32) subject to (7.33) and (7.34), that is, it reaches the maximum
        value $V(t_0, x(t_0))$. Then
        \begin{align}
            V(t_0, x(t_0)) &= \int_{t_0}^{t_1} f(t, \hat{x}(t), \hat{y}(t))dt
            + V(t_1, \hat{x}(t_1)) \tag{7.38}\\
            \nonumber &= \max_{y(t) \in \Y} \left\{ \int_{t_0}^{t_1} f(t,
                x(t), y(t))dt + V(t_1, x(t_1)) \right\}
        \end{align}
        for all $t_1 \geq t_0$, where in the second equation the trajectory
        of $x(t)$ is given by $\dot x(t) = g(t, x(t), y(t))$.
    \end{block}
    \begin{itemize}
        \item Analogous to the Principle of Optimality in dynamic programming
        \item Discounting is embedded in $f$
    \end{itemize}
\end{frame}

\begin{frame}{Infinite-Horizon Maximum Principle}
\begin{block}{Theorem 7.9 (Infinite-Horizon Maximum Principle)}
    Suppose that the problem of maximizing (7.32) subject to (7.33) and
    (7.34), with $f$ and $g$ continuously differentiable, has a piecewise
    continuous interior solution $(\hat{x}(t), \hat{y}(t)) \in \text{Int } \X
    \times \text{Int } \Y$. Let $H(t, x(t), y(t), \lambda(t)) \equiv f(t, x(t), y(t)) +
    \lambda(t)g(t, x(t), y(t))$. Then given $(\hat{x}(t), \hat{y}(t))$, the
    Hamiltonian $H(t, x, y, \lambda)$ satisfies the Maximum Principle:
    \begin{align*}
        H(t, \hat{x}(t), \hat{y}(t), \lambda(t)) \geq H(t, \hat{x}(t), y(t), \lambda(t))
    \end{align*}
    for all $y(t) \in \Y$ and for all $t \in \mathbb{R}_+$. Moreover, for all
    $t \in \mathbb{R}_+$ for which $\hat{y}(t)$ is continuous, the following
    necessary conditions are satisfied:
    \begin{align}
        &H_y(t, \hat{x}(t), \hat{y}(t), \lambda(t)) = 0, \tag{7.39} \\
        &\dot{\lambda}(t) = -H_x(t, \hat{x}(t), \hat{y}(t), \lambda(t)),
        \tag{7.40}
    \end{align}
    and
    \begin{align}
        &\dot{x}(t) = H_\lambda(t, \hat{x}(t), \hat{y}(t), \lambda(t)),
        \text{ with } x(0) = x_0 \text{ and } \lim_{t \to \infty} b(t)x(t)
        \geq x_1. \tag{7.41} 
    \end{align}

\end{block}
\end{frame}

\begin{frame}{Remarks on Theorem 7.9}
    \begin{itemize}
        \item Theorem 7.9 can be viewed as stronger than the theorems
        presented in the discrete time case, especially since it does not
        impose compactness-type conditions.
        \item Nevertheless, this theorem only applies when the maximization problem
        has a piecewise continuous solution $\hat{y}(t)$.
        \item Economic problems often have enough structure to ensure that
        $\hat{y}(t)$ is indeed a continuous function of $t$. Consequently, in
        most economic problems it is sufficient to focus on the necessary
        conditions (7.39)–(7.41).
\end{itemize}
\end{frame}


\begin{frame}{Hamilton-Jacobi-Bellman Equation}
\begin{block}{Theorem 7.10 (Hamilton-Jacobi-Bellman Equation)}
    Let $V(t, x)$ be as defined in (7.35), and suppose that the hypotheses in
    Theorem 7.9 hold. Then when $V(t, x)$ is differentiable in $(t, x)$, $V$
    satisfies the HJB equation
    \begin{equation*}
        -\frac{\partial V(t, x)}{\partial t} = \max_{y \in \Y} \left\{ f(t, x,
              y) + \frac{\partial V(t, x)}{\partial x} g(t, x, y) \right\}.
    \end{equation*}
    The optimal pair $(\hat{x}(t), \hat{y}(t))$ satisfies:
    \begin{align*}
        -\frac{\partial V(t, \hat{x}(t))}{\partial t} &= f(t, \hat{x}(t),
        \hat{y}(t)) + \frac{\partial V(t, \hat{x}(t))}{\partial x} g(t,
        \hat{x}(t), \hat{y}(t)) \\
        &= \max_{y \in \Y} \left\{f(t, \hat{x}(t), y) + \frac{\partial V(t,
              \hat{x}(t))}{\partial x} g(t, \hat{x}(t), y) \right\}
    \end{align*}
    for all $t \in \mathbb{R}_+$.
\end{block}
\textbf{Remark.} The HJB equation may admit multiple solutions. An appropriate
asymptotic, growth, or transversality condition is generally needed to select
the value-function solution. If the discounted continuation value vanishes,
this condition may be $\lim_{t\to\infty}V(t,x)=0$.
\end{frame}

\begin{frame}{Importance and Features of HJB Equation}
\begin{itemize}
    \item The HJB equation is a partial differential equation that is useful
    for providing an intuition for the Maximum Principle.
    \item This partial differential equation also has a similarity to the
    Euler equation derived in the context of discrete-time dynamic
    programming: the first term on the right-hand side corresponds to the
    current gain and the second term to the benefit of increasing the state.
    \item The left-hand side results from the fact that the maximized value
    can also change over time.
    \item The HJB equation implies that current gain is equal to the loss of
    value over time:
    \begin{equation*}
        f(t, \hat x(t), \hat y(t)) = -\frac{d}{dt}V(t, \hat x(t))
    \end{equation*}
\end{itemize}
\end{frame}

% \begin{frame}{Sufficiency Conditions}
% \begin{block}{Theorem 7.11 (Sufficiency Conditions for Infinite-Horizon Optimal Control)}
% Consider the problem of maximizing (7.32) subject to (7.33) and (7.34), with $f$ and $g$ continuously differentiable. Define $H(t, x, y, \lambda)$ as in (7.16), and suppose that an admissible pair $(\hat{x}(t), \hat{y}(t))$ satisfies (7.39)–(7.41).

% Given the resulting costate variable $\lambda(t)$, define $M(t, x, \lambda) \equiv \max_{y(t) \in \Y(t)} H(t, x, y, \lambda)$ as in (7.20).

% If $\X$ is a convex set, $M(t, x, \lambda)$ is concave in $x \in \X$ for all $t$ and $\lim_{t \to \infty} \lambda(t)(\hat{x}(t) - \tilde{x}(t)) \leq 0$ for all $\tilde{x}(t)$ implied by an admissible control path $\tilde{y}(t)$, then the pair $(\hat{x}(t), \hat{y}(t))$ achieves the global maximum of (7.32).

% Moreover, if $M(t, x, \lambda)$ is strictly concave in $x$, then $(\hat{x}(t), \hat{y}(t))$ is the unique solution to (7.32).
% \end{block}
% \end{frame}

% \begin{frame}{Remarks on Sufficiency}
% This sufficiency result involves the difficult-to-check condition $\lim_{t \to \infty} \lambda(t)(x(t) - \tilde{x}(t)) \leq 0$ for all $\tilde{x}(t)$ (i.e., for all $\tilde{x}(t)$ implied by an admissible control path $\tilde{y}(t)$).

% When we impose the appropriate transversality condition, sufficiency becomes much more straightforward to verify.
% \end{frame}

\begin{frame}
    \frametitle{Heuristic Derivation of the Maximum Principle}
    We can use the HJB equation to derive the Maximum Principle by setting
    the costate to
        \begin{equation*}
            \lambda(t) = \frac{\partial V(t, \hat{x}(t))}{\partial x}
        \end{equation*}
    \begin{itemize}
        \item Then the first order condition for the HJB equation gives
        $H_y(t, \hat{x}(t), \hat{y}(t), \lambda(t)) = 0$
        \item Taking the derivative of both sides of the HJB equation with
        respect to $x$ gives
        \begin{equation*}
            - \frac{\partial^2 V(t, \hat x(t))}{\partial t \partial x} =
            f_x(t, \hat{x}(t), \hat{y}(t)) + \frac{\partial^2 V(t,
              \hat{x}(t))}{\partial x^2} g(t, \hat{x}(t), \hat{y}(t)) +
            \frac{\partial V(t, \hat{x}(t))}{\partial x} g_x(t, \hat{x}(t),
            \hat{y}(t))
        \end{equation*}
        \item By the definition of $\lambda(t)$, we have
        \begin{equation*}
            \dot \lambda(t) = \frac{\partial^2 V(t, \hat x(t))}{\partial t
              \partial x} + \frac{\partial^2 V(t, \hat{x}(t))}{\partial x^2}
            g(t, \hat{x}(t), \hat{y}(t))
        \end{equation*}
        \item Combining the above equations gives the costate equation
        \begin{equation*}
            \dot{\lambda}(t) = -H_x(t, \hat{x}(t), \hat{y}(t), \lambda(t))
        \end{equation*}
    \end{itemize}
\end{frame}

\begin{frame}{Economic Intuition}
    \begin{itemize}
        \item Since $\lambda(t) = \frac{\partial V(t, \hat{x}(t))}{\partial
          x}$, $\lambda(t)$ measures the impact (shadow value) of a small
        increase in $x$ on the optimal value of the program.
        \item Consider the problem of maximizing ($L$ can be thought of as
        the Lagrangian)
        \begin{align*}
            \int_0^{t_1} L(t, \hat x(t), y(t), \lambda(t)) &\equiv \int_0^{t_1}
            \left[f(t, \hat{x}(t), y(t)) + \lambda(t)\left( g(t,
                    \hat{x}(t), y(t)) - \dot x(t) \right)\right]dt\\
            &= \int_0^{t_1} \left[H(t, \hat{x}(t), y(t), \lambda(t)) - \lambda(t)
            \dot x(t)\right] dt
        \end{align*}
        with respect to the entire function $y(t)$. The necessary condition
        is $H_y(t, \hat{x}(t), y(t), \lambda(t)) = 0$.
        \item The maximum principle implies that $f_y(t, \hat x(t), \hat
        y(t)) + \lambda(t) g_y(t, \hat x(t), \hat y(t)) = 0$: the marginal
        effect of a change in $y(t)$ on instantaneous payoff should
        counter-balance that on the value of stock.
        % \item The problem is equivalent to maximizing instantaneous returns,
        % plus the value of stock of $x(t)$, as given by $\lambda(t)$, times
        % the increase in the stock.
    \end{itemize}
\end{frame}

\begin{frame}{Costate Equation Interpretation}
\begin{itemize}
    \item In order for the above explanation to make sense, $\lambda(t)$ must
    follow the costate equation,
    \begin{align*}
        -\dot{\lambda}(t) &= H_x(t, \hat{x}(t), \hat{y}(t), \lambda(t)) \\
        &= f_x(t, \hat{x}(t), \hat{y}(t)) + \lambda(t)g_x(t, \hat{x}(t), \hat{y}(t))
    \end{align*}
    \item $\dot{\lambda}(t)$ is the appreciation rate in the stock variable
    $x(t)$
    \item A small increase in the state $x(t)$ changes the current flow
    return by $f_x(t, \hat{x}(t), \hat{y}(t))$ and also changes the value of
    the stock by $\lambda(t)g_x(t, \hat{x}(t), \hat{y}(t))$
    \item This gain should be equal to the depreciation in the value of the
    stock $-\dot{\lambda}(t)$ over time
\end{itemize}
\end{frame}


\begin{frame}{Stationary HJB}
    Given its prominent role in dynamic economic analysis, it is useful to
    consider the simpler stationary version of the HJB equation:
 \begin{itemize}
        \item $f(t, x(t), y(t)) = \exp(-\rho t)f(x(t), y(t))$
        \item The law of motion of the state variable is given by an
        autonomous differential equation, that is, $g(t, x(t), y(t)) =
        g(x(t), y(t))$.
    \end{itemize}

    In this case, one can easily verify that if an admissible pair
    $(\hat{x}(t), \hat{y}(t))_{t \geq 0}$ is optimal starting at $t = 0$ with
    initial condition $x(0) = x_0$, then its continuation from any $s>0$ is
    optimal starting from the reached state $x(s)=\hat{x}(s)$.
\end{frame}

\begin{frame}{Stationary HJB}
    \begin{itemize}
        \item Let us define $v(x) \equiv V(0, x)$. Since $(\hat{x}(t),
        \hat{y}(t))$ is an optimal plan regardless of the starting date, we
        have
        \begin{align}
            V(t, x(t)) = \exp(-\rho t)v(x(t)) \text{ for all } t. \tag{7.43}
        \end{align}
        \item Then by definition,
        \begin{align*}
            \frac{\partial V(t, x(t))}{\partial t} = -\rho \exp(-\rho
            t)v(x(t)).
        \end{align*}
        \item Then we can derive the stationary form of the HJB equation
        \begin{align}
            \rho v(\hat{x}(t)) = f(\hat{x}(t), \hat{y}(t)) +
            \dot{v}(\hat{x}(t)). \tag{7.44}
        \end{align}
    \end{itemize}
    This stationary HJB equation is widely used in dynamic economic analysis
    and can be interpreted as a ``no-arbitrage asset value equation''.
\end{frame}

\begin{frame}{Economic Intuition of the Stationary HJB Equation}
    \begin{itemize}
        \item The stationary form of the HJB equation is
        \begin{align}
            \rho v(\hat{x}(t)) = f(\hat{x}(t), \hat{y}(t)) +
            \dot{v}(\hat{x}(t)), \tag{7.48} 
        \end{align}
        \item We can think of $v$ as the value of an asset and $\rho$ as the
        required rate of return
        \begin{itemize}
            \item Dividends are given by the flow payoff $f(\hat{x}(t), \hat{y}(t))$
            \item Capital gains or losses are given by $\dot{v}$
        \end{itemize}
        \item In equilibrium, the return on this asset is equal to the
        required rate of return $\rho$.
        \item The Maximum Principle (for stationary problems) requires that
        $v(x)$, and its rate of change, $\dot{v}(x)$, should be consistent
        with this no-arbitrage condition.
    \end{itemize}
\end{frame}

\section{Discounted Infinite-Horizon Optimal Control}

\begin{frame}
\frametitle{Problem Formulation}
\begin{itemize}
    \item Economically interesting problems often take a more specific form
    with exponential discounting:
    \begin{align}
        \max_{x(t),y(t)} W(x(t), y(t)) &\equiv \int_0^\infty \exp(-\rho
        t)f(x(t), y(t))dt, \quad \text{with } \rho > 0 \tag{7.60} 
    \end{align}
    subject to
    \begin{align}
        \dot{x}(t) &= g(t, x(t), y(t)) \tag{7.61}
    \end{align}
    and
    \begin{align}
        x(t) &\in \text{Int } \X(t) \text{ and } y(t) \in \text{Int } \Y(t) \text{ for all } t, \nonumber \\
        x(0) &= x_0, \text{ and } \lim_{t\to\infty} b(t)x(t) \geq x_1 \tag{7.62}
    \end{align}
    \item Assume positive discounting: $\rho > 0$
    \item Recall that $b : \mathbb{R}_+ \to \mathbb{R}_+$ and
    $\lim_{t\to\infty} b(t) < \infty$
\end{itemize}
\end{frame}

\begin{frame}
    \frametitle{Current-Value Hamiltonian}
    \begin{itemize}
        \item The Hamiltonian in this case is:
        \begin{align*}
            H(t, x(t), y(t), \lambda(t)) &= \exp(-\rho t)f(x(t), y(t)) +
            \lambda(t)g(t, x(t), y(t)) \\ 
            &= \exp(-\rho t)[f(x(t), y(t)) + \mu(t)g(t, x(t), y(t))]
        \end{align*}
        where the second line uses the definition
        \begin{align}
            \mu(t) \equiv \exp(\rho t)\lambda(t) \tag{7.63}
        \end{align}
        \item We can work with the current-value Hamiltonian, defined as:
        \begin{align}
            \hat{H}(t, x(t), y(t), \mu(t)) \equiv f(x(t), y(t)) + \mu(t)g(t,
            x(t), y(t)) \tag{7.64} 
        \end{align}
        \item When $g(t, x(t), y(t))$ is also an autonomous differential
        equation of the form $g(x(t), y(t))$, we can simply write
        $\hat{H}(x(t), y(t), \mu(t))$ 
    \end{itemize}
\end{frame}

\begin{frame}
\frametitle{Assumptions}

Throughout, $f$ and $g$ are continuously differentiable for all admissible
$(x(t), y(t))$, with derivatives denoted $f_x$, $f_y$, $g_x$, and $g_y$.
\begin{block}{Assumption 7.1}
    In the maximization of (7.60) subject to (7.61) and (7.62):
    \begin{enumerate}
        \item $f$ is weakly monotone in $x$ and $y$, and $g$ is weakly
        monotone in $(t, x, y)$ 
        \item there exists $m > 0$ such that $|g_y(t, x(t), y(t))| \geq m$
        for all $t$ and for all admissible pairs $(x(t), y(t))$
        \item there exists $M < \infty$ such that $|f_y(x, y)| \leq M$ for
        all $x$ and $y$
    \end{enumerate}
\end{block}
\end{frame}


\begin{frame}
\frametitle{Maximum Principle}
\begin{block}{Theorem 7.13 (Maximum Principle for Discounted Infinite-Horizon Problems)}
    Suppose that the problem of maximizing (7.60) subject to (7.61) and
    (7.62), with $f$ and $g$ continuously differentiable, has an interior
    piecewise continuous optimal control $\hat{y}(t) \in \text{Int } \Y(t)$
    with corresponding state variable $\hat{x}(t) \in \text{Int } \X(t)$.
    Suppose that the value function $V(t, \hat{x}(t))$ is differentiable in
    $x$ and $t$ for $t$ sufficiently large, that $V(t, \hat{x}(t))$ exists
    and is finite for all $t$, and that $\lim_{t\to\infty} \partial V(t,
    \hat{x}(t))/\partial t = 0$. Then except at points of discontinuity of
    $\hat{y}(t)$, the optimal control pair $(\hat{x}(t), \hat{y}(t))$
    satisfies the following necessary conditions:
    \begin{gather}
        \hat{H}_y(t, \hat{x}(t), \hat{y}(t), \mu(t)) = 0 \text{ for all } t
        \in \mathbb{R}_+ \tag{7.65} \\
        \rho\mu(t) - \dot{\mu}(t) = \hat{H}_x(t, \hat{x}(t), \hat{y}(t),
        \mu(t)) \text{ for all } t \in \mathbb{R}_+ \tag{7.66}\\
        \dot{x}(t) = \hat{H}_\mu(t, \hat{x}(t), \hat{y}(t), \mu(t)) \text{
          for all } t \in \mathbb{R}_+, x(0) = x_0, \text{ and }
        \lim_{t\to\infty} b(t)x(t) \geq x_1 \tag{7.67}
    \end{gather}
\end{block}
\end{frame}

\begin{frame}
\frametitle{Maximum Principle}
\begin{block}{Theorem 7.13 (Continued)}
    and the transversality condition
    \begin{align}
        \lim_{t\to\infty}[\exp(-\rho t)\hat{H}(t, \hat{x}(t), \hat{y}(t),
        \mu(t))] = 0 \tag{7.68} 
    \end{align}
    Moreover, suppose that Assumption 7.1 holds and that either
    $\lim_{t\to\infty} \hat{x}(t) = x^* \in \mathbb{R}$ or $\lim_{t\to\infty}
    \dot{x}(t)/\hat{x}(t) = \chi \in \mathbb{R}$. Then the transversality
    condition can be strengthened to
    \begin{align}
        \lim_{t\to\infty}[\exp(-\rho t)\mu(t)\hat{x}(t)] = 0 \tag{7.69}
    \end{align}
\end{block}

\end{frame}

\begin{frame}
\frametitle{The Transversality Condition}
\begin{itemize}
    \item Notice that compared to the transversality condition in the
    finite-horizon case, there is the additional term $\exp(-\rho t)$ in
    (7.69). This is because the transversality condition applies to the
    original costate variable $\lambda(t)$:
    $\lim_{t\to\infty}[\lambda(t)x(t)] = 0$
    \item Note also that the stronger transversality condition takes the form
    \begin{align*}
        \lim_{t\to\infty}[\exp(-\rho t)\mu(t)\hat{x}(t)] = 0
    \end{align*}
    not simply $\lim_{t\to\infty}[\exp(-\rho t)\mu(t)] = 0$.
\end{itemize}
\end{frame}

\begin{frame}
\frametitle{Limitations}
\begin{itemize}
    \item It is important to emphasize that Theorem 7.13 only provides
    necessary conditions for interior continuous solutions (with
    $\lim_{t\to\infty} \hat{x}(t) = x^*$ or $\lim_{t\to\infty}
    \dot{x}(t)/\hat{x}(t) = \chi$).
    \item The next theorem shows that under the appropriate concavity
    conditions, (7.69) is also a sufficient transversality condition.
    \item It further shows that for such concave problems, Assumption 7.1 or
    the limiting conditions in Theorem 7.13 are no longer required.
\end{itemize}
\end{frame}

\begin{frame}
\frametitle{Sufficiency Conditions}
\begin{block}{Theorem 7.14 (Sufficiency Conditions for Discounted
      Infinite-Horizon Problems)}
    Consider the problem of maximizing (7.60) subject to (7.61) and (7.62),
    with $f$ and $g$ continuously differentiable. Suppose that some
    $\hat{y}(t)$ and the corresponding path of state variable $\hat{x}(t)$
    satisfy the necessary conditions and the transversality condition
    (7.65)--(7.68). Given the resulting current-value costate variable
    $\mu(t)$, define $M(t, x, \mu) \equiv \max_{y(t)\in \Y(t)} \hat{H}(t, x,
    y, \mu)$. Suppose that
    \begin{itemize}
        \item $V(t, \hat{x}(t))$ exists and is finite for all $t$;
        \item for any admissible pair $(x(t), y(t))$,
        $\lim_{t\to\infty}[\exp(-\rho t)\mu(t)x(t)] \geq 0$;
        \item $\X(t)$ is convex and $M(t, x, \mu)$ is concave in $x \in
        \X(t)$ for all $t$.
    \end{itemize}
    Then the pair $(\hat{x}(t), \hat{y}(t))$ achieves the global maximum of
    (7.60). Moreover, if $M(t, x, \mu)$ is strictly concave in $x$,
    $(\hat{x}(t), \hat{y}(t))$ is the unique solution to (7.60).
\end{block}

\end{frame}
\begin{frame}
\frametitle{Sufficiency Conditions}

Theorem 7.14 is very useful and powerful. Given this result, the following
strategy will be used in most problems:
\begin{enumerate}
    \item Use the conditions in Theorem 7.13 to locate a candidate interior
    solution $(\hat{x}(t), \hat{y}(t))$ satisfying (7.65)--(7.68)
    \item Then verify the concavity conditions of Theorem 7.14 and simply
    check that $\lim_{t\to\infty}[\exp(-\rho t)\mu(t)x(t)] \geq 0$ for other
    admissible pairs, with $\mu(t)$ associated with the candidate solution
    $(\hat{x}(t), \hat{y}(t))$
    \item If these conditions are satisfied, we will have characterized a
    global maximum
\end{enumerate}
\end{frame}

\begin{frame}
\frametitle{Corollary 7.1: Continuity}
\begin{block}{Corollary 7.1}
    Suppose that the hypotheses in Theorem 7.14 are satisfied, $M(t, x, \mu)$
    is strictly concave in $x$ for all $t$, and $\Y$ is compact. Then
    $\hat{y}(t)$ is a continuous function of $t$ on $\mathbb{R}_+$.
\end{block}
\end{frame}

\end{document}
